\section*{Capítulo \ref{ch:g12:ka-and-pka} --- Ácidos y bases: Ka y pKa}

\begin{solution}{exo:g12:ka-and-pka:1}
Bases conjugadas: \ce{HCOO-}, \ce{NH3}, \ce{OH-}, \ce{CO3^{2-}}. Ácidos
conjugados: \ce{NH4+}, \ce{H2O}, \ce{H3O+}, \ce{CO2,H2O}. \omterm{def:g12:ka-and-pka:oxonium}{Anfolitos}:
\ce{H2O} y \ce{HCO3-}.
\end{solution}

\begin{solution}{exo:g12:ka-and-pka:2}
\omterm{def:g9:acids-bases-ph:ph}{pH} 1.0; 3.0; $-\log(\num{5.0e-3}) = 2.30$.
\end{solution}

\begin{solution}{exo:g12:ka-and-pka:3}
$10^{-4.5} = \qty{3.2e-5}{mol/L}$. A \omterm{def:g9:acids-bases-ph:ph}{pH} 11.2:
$[\ce{H3O+}] = \qty{6.3e-12}{mol/L}$ y
$[\ce{OH-}] = \num{1.0e-14}/\num{6.3e-12} = \qty{1.6e-3}{mol/L}$.
\end{solution}

\begin{solution}{exo:g12:ka-and-pka:4}
\ce{HCOOH + H2O <=> HCOO- + H3O+},
$K_a = \dfrac{[\ce{HCOO-}][\ce{H3O+}]}{[\ce{HCOOH}]\,c^\circ}$;
\ce{NH4+ + H2O <=> NH3 + H3O+},
$K_a = \dfrac{[\ce{NH3}][\ce{H3O+}]}{[\ce{NH4+}]\,c^\circ}$.
\end{solution}

\begin{solution}{exo:g12:ka-and-pka:5}
El ácido metanoico (menor $pK_a$). La base más fuerte es \ce{CH3COO-}, la
base del ácido más débil.
\end{solution}

\begin{solution}{exo:g12:ka-and-pka:6}
$[\ce{A-}] = [\ce{H3O+}] = \qty{1.0e-3}{mol/L}$;
$[\ce{HA}] = 0.050 - 0.001 = \qty{0.049}{mol/L}$;
$K_a = \num{1.0e-6}/0.049 = \num{2.0e-5}$; $pK_a = 4.69$.
\end{solution}

\begin{solution}{exo:g12:ka-and-pka:7}
$K_a = \num{6.31e-5}$; $x^2 + \num{6.31e-5}\,x - \num{1.26e-6} = 0$ da
$x = \qty{1.09e-3}{mol/L}$: \omterm{def:g9:acids-bases-ph:ph}{pH} 2.96.
\end{solution}

\begin{solution}{exo:g12:ka-and-pka:8}
$\mathrm{pH} = 14.0 + \log(\num{2.0e-3}) = 14.0 - 2.70 = 11.3$.
\end{solution}

\begin{solution}{exo:g12:ka-and-pka:9}
\ce{HF}, \ce{HCOOH}, ácido ascórbico, \ce{C6H5COOH}, \ce{CH3COOH},
\ce{CO2,H2O}, \ce{NH4+}, \ce{HCO3-}. Un \omterm{def:g12:ka-and-pka:strong-acid}{ácido fuerte} se sitúa muy abajo,
más allá del extremo de la escala: reacciona totalmente con el agua.
\end{solution}

\begin{solution}{exo:g12:ka-and-pka:10}
$[\ce{H3O+}] = 10^{-5.6} = \qty{2.5e-6}{mol/L}$, 25 veces más que en el
agua pura (\qty{1.0e-7}{mol/L}).
\end{solution}

\begin{solution}{exo:g12:ka-and-pka:11}
$K_a K_b = \dfrac{[\ce{NH3}][\ce{H3O+}]}{[\ce{NH4+}]c^\circ} \times
\dfrac{[\ce{NH4+}][\ce{OH-}]}{[\ce{NH3}]c^\circ} = K_e$, así que
$pK_a = pK_e - pK_b = 14.0 - 4.74 = 9.26$.
\end{solution}

\begin{solution}{exo:g12:ka-and-pka:12}
A \qty{0.0010}{mol/L}: $x^2 + \num{1.74e-5}\,x - \num{1.74e-8} = 0$,
$x = \qty{1.24e-4}{mol/L}$, \omterm{def:g9:acids-bases-ph:ph}{pH} 3.91: una subida de 0.52, no de 1. Al
diluirse, un \omterm{def:g12:ka-and-pka:strong-acid}{ácido débil} reacciona en mayor proporción con el agua (aquí
el \qty{12}{\percent} frente al \qty{4}{\percent}), lo que compensa en
parte la \omterm{def:g10:concentration-and-dilution:dilution}{dilución}.
\end{solution}

\begin{solution}{exo:g12:ka-and-pka:13}
Como ácido: \ce{HCO3- + H2O <=> CO3^{2-} + H3O+}, $K = 10^{-10.33}$.
Como base: \ce{HCO3- + H2O <=> H2CO3 + OH-} (donde \ce{H2CO3} representa \ce{CO2,H2O}), $K = K_e/K_a = 10^{-14.0+6.37}
= 10^{-7.63}$. La reacción como base tiene la constante mayor: la
disolución es ligeramente básica.
\end{solution}

\begin{solution}{exo:g12:ka-and-pka:14}
Un ácido, por diluido que esté, no puede volver básica una disolución.
Cuando los \omterm{def:g12:ka-and-pka:oxonium}{iones oxonio} propios del ácido son menos que los
\qty{1e-7}{mol/L} que aporta el agua, dominan los \omterm{def:g9:ions:ion}{iones} del agua: el \omterm{def:g9:acids-bases-ph:ph}{pH}
se acerca a 7 por debajo y nunca lo supera.
\end{solution}

\begin{solution}{exo:g12:ka-and-pka:15}
$n = 0.500 / 176.0 = \qty{2.84e-3}{mol}$, $c = \qty{0.0142}{mol/L}$;
$K_a = \num{9.12e-5}$; $x^2 + \num{9.12e-5}\,x - \num{1.30e-6} = 0$ da
$x = \qty{1.09e-3}{mol/L}$: \omterm{def:g9:acids-bases-ph:ph}{pH} 2.96.
\end{solution}

\begin{solution}{pb:g12:ka-and-pka:1}
\textbf{1.} \ce{H-C(=O)-O-H}: el hidrógeno que cede es el del \ce{O-H}.

\textbf{2.} \ce{HCOOH}/\ce{HCOO-}; \ce{HCOOH + H2O <=> HCOO- + H3O+}.

\textbf{3.} $K_a = \dfrac{[\ce{HCOO-}][\ce{H3O+}]}{[\ce{HCOOH}]c^\circ} =
10^{-3.74} = \num{1.82e-4}$.

\textbf{4.} Por debajo del ácido etanoico en la escala ($3.74 < 4.76$):
más fuerte.

\textbf{5.} $M = \qty{46.0}{g/mol}$: $0.010 \times 46.0 = \qty{0.46}{g}$.

\textbf{6.} \ce{HCOOH}: $0.010 - x$; \ce{HCOO-}: $x$; \ce{H3O+}: $x$.

\textbf{7.} $\dfrac{x^2}{0.010 - x} = \num{1.82e-4}$.

\textbf{8.} $x^2 + \num{1.82e-4}\,x - \num{1.82e-6} = 0$:
$x = \qty{1.26e-3}{mol/L}$.

\textbf{9.} $\mathrm{pH} = -\log(\num{1.26e-3}) = 2.90$.

\textbf{10.} $\num{1.26e-3}/0.010 = \qty{12.6}{\percent}$.

\textbf{11.} \qty{1.26e-3}{mol/L} es más de diez mil veces
\qty{1e-7}{mol/L}: los \omterm{def:g9:ions:ion}{iones} propios del agua son despreciables.

\textbf{12.} \omterm{def:g9:acids-bases-ph:ph}{pH} 2.0.

\textbf{13.} $0.010 / \num{1.26e-3} \approx 8$ veces más.

\textbf{14.} El ácido clorhídrico es un \omterm{def:g12:ka-and-pka:strong-acid}{ácido fuerte}, transformado por
completo en \omterm{def:g9:ions:ion}{iones}; el ácido metanoico es débil: su reacción con el agua
se detiene en un equilibrio en el que la mayor parte queda como
\ce{HCOOH}.

\textbf{15.} \ce{HCOOH + HCO3- -> HCOO- + CO2 + H2O}.

\textbf{16.} Multiplicando el numerador y el denominador de
\[
  K = \frac{[\ce{HCOO-}]\,[\ce{CO2}]}{[\ce{HCOOH}]\,[\ce{HCO3-}]}
\]
por $[\ce{H3O+}]$ se obtiene $K = K_a(\ce{HCOOH}) / K_a(\ce{CO2,H2O}) =
\num{1.82e-4}/\num{4.3e-7} \approx 420$.

\textbf{17.} Sí, $K \gg 1$; se desprenden burbujas de dióxido de carbono.

\textbf{18.} El hidróxido es una \omterm{def:g12:ka-and-pka:strong-acid}{base fuerte} y atacaría por sí mismo la
piel; el hidrogenocarbonato es una \omterm{def:g12:ka-and-pka:strong-acid}{base débil} que neutraliza el ácido y
resulta suave aunque esté en exceso.

\textbf{19.} \omterm{def:g9:acids-bases-ph:ph}{pH} 2.9.
\end{solution}
